\section{Conclusion}\label{sec:conclusion_main}
\noindent This report presented the design, implementation, and empirical evaluation of the Lightweight Blockchain-Based Data Trading (LBDT) framework for the Internet of Vehicles (IoV). LBDT successfully resolves the critical conflict between transaction confirmation latency, consensus message overhead, node reputation evaluation, and fair market pricing in high-mobility vehicular environments.

By decoupling the blockchain ledger into Keyblocks ($\tau_{key} = 60\text{ s}$) and Microblocks ($\tau_{micro} = 6\text{ s}$), LBDT achieves rapid local transaction settlement while preserving global state security. The model of aged Gompertz reputation, which is powered by a trained logistic regression classifier (phonest ROC-AUC = 0.9991), is used dynamically to detect and eradicate malicious nodes that are involved in fraudulent trades or dropping packets.

Furthermore, the mechanism of twofold auction based on reputation penalizes ask prices of sellers based on delays in data generation and reputation deficiencies of sellers. 
The consensus pipeline, which uses EC-VRF leader election with a 4-phase Gosig BFT protocol, limits the complexity of messages to $5N-3$ messages within a round and keeps the safety of the committee with a ratio of malicious reputation equal to 19.4\% (len under BFT limit of 33.3\%).Numerous simulations of 1,000 active vehicles and 400 road-side units were carried out using SUMO and Veins, confirming that the average time of transaction confirmation made 2.4828 seconds and proving that LBDT can be used for practical purposes in the real-time application for vehicles.

\section{Summary of Key Contributions}\label{sec:conclusion_contributions}
\noindent The key outputs of this project include:
\begin{enumerate}
    \item \textbf{Decoupled Dual-Chain Ledger:} Which allows achieving sub-three-second transaction confirmation through local microblock processing taking six seconds and separating it from the 60-second epoch control of the keyblock.
    \item \textbf{ML-Enhanced Gompertz Trust Model:} Where the exponential decline of aging in the human population has been combined with predictions in Logistic Regression to secure defense against oscillating and whitewashing attacks.
    \item \textbf{Fair Double Auction Engine:} where the concept of penalties for unsuccessful bids ($loss = \alpha_1 (\Delta t)^{\beta_1} + \alpha_2 (1 - Rep_i)^{\beta_2}$) is utilized to prevent data purchasers from receiving bad quality data or data with high latency.
    \item \textbf{Lightweight Gosig BFT Consensus:} leading to a reduction in wireless bandwidth usage by over 78\% through limiting BFT communication to $5N-3$ messages.
    \item \textbf{Dynamic Load-Aware Sharding:} Allowing for automated splitting and merging of geographic clusters based on live zone utility scores ($U_z$).
\end{enumerate}

\section{System Limitations}\label{sec:conclusion_limitations}
\noindent Though these positive outcomes seem to be promising, some issues can be identified in the current execution:
\begin{itemize}
    \item \textbf{Cryptographic Verification Latency on OBUs:}Although RSUs conduct BFT concurrence smoothly, OBUs in vehicles with limited resources deal with insignificant load on the CPU when validating goatsig signatures in peak periods.
    \item \textbf{Private Bid Disclosure:} Double auction quotations are accessible in the RSUs of the zone. Although hashes of transaction data have been protected, there is still a risk of revealing bidding prices and thus heating insights into marketing strategy.
\end{itemize}

\section{Future Scope}\label{sec:conclusion_future}
\noindent Future research approaches to advance and improve the LBDT framework comprise the following: 
\begin{enumerate}
    \item \textbf{Zero-Knowledge Proofs (ZK-SNARKs) for Bid Privacy:} Adding ZKSNARK into bidding in double auctions to allow for verifiable price matching without disclosing raw bid amounts to zone RSUs. 
    \item \textbf{Deep Reinforcement Learning (DRL) for Sharding:} Using DRL agents instead of the heuristic utility score ($U_z$) for predicting the traffic congestion and obtaining proactive zone sharding recommendations. 
    \item \textbf{5G/6G V2X Slicing Integration:} Using the 5G/6G network slicing for assigning reliable URLLC channels to the microblock consensus data transfers.
    \item \textbf{Hardware Security Module (HSM) Deployment:} Integrating OBU Hardware Security Modules (e.g., TPM 2.0) for tamper-proof key storage and cryptographic signature acceleration.
\end{enumerate}
